\documentclass[11pt, a4paper]{article}

% ── Packages ──────────────────────────────────────────────────────────
\usepackage[utf8]{inputenc}
\usepackage[T1]{fontenc}
\usepackage{lmodern}
\usepackage[margin=1in]{geometry}
\usepackage{amsmath, amssymb, amsthm}
\usepackage{booktabs}
\usepackage{longtable}
\usepackage{graphicx}
\usepackage{hyperref}
\usepackage{xcolor}
\usepackage{listings}
\usepackage{tcolorbox}
\usepackage{polya}

\hypersetup{
  colorlinks=true,
  linkcolor=polyablue,
  urlcolor=polyablue,
  citecolor=polyablue,
}

% ── Code listing style ────────────────────────────────────────────────
\lstdefinestyle{polya-python}{
  language=Python,
  basicstyle=\ttfamily\small,
  keywordstyle=\color{polyablue}\bfseries,
  commentstyle=\color{polyagray}\itshape,
  stringstyle=\color{polyagreen},
  numbers=left,
  numberstyle=\tiny\color{polyagray},
  numbersep=8pt,
  frame=single,
  framerule=0.4pt,
  rulecolor=\color{polyagray},
  breaklines=true,
  showstringspaces=false,
  tabsize=4,
}
\lstset{style=polya-python}
\title{Anomaly Detection -- Server Response Times}
\author{uber-polya}
\date{2026-02-26}

\begin{document}

\maketitle
\tableofcontents
\newpage

% ── Phase A: Problem Formulation ─────────────────────────────────────
\section{Problem Formulation}

\begin{polyamodel}

\textbf{Problem Type}: Find \hfill
\textbf{Domain}: Time Series Analysis


\end{polyamodel}

% ── Universe ──────────────────────────────────────────────────────────
\subsection{Universe}

\begin{itemize}
  \item $Zscore Detection: \{n_detected, indices, precision, recall, f1\}$
  \item $Iqr Detection: \{n_detected, indices, precision, recall, f1\}$
  \item $Combined Detection: \{n_detected, indices, precision, recall, f1\}$
  \item $Change Point: \{detected, best_match, true, error\}$
  \item $Regime Analysis: \{pre_change_mean, post_change_mean, t_statistic, p_value\}$
\end{itemize}

% ── Variables ─────────────────────────────────────────────────────────
\subsection{Variables}

\begin{longtable}{lll}
\toprule
\textbf{Symbol} & \textbf{Meaning} & \textbf{Type / Range} \\
\midrule
\endhead
$x$ & decision variable & $\mathbb{{R}}$ \\
\bottomrule
\end{longtable}

% ── Structure ─────────────────────────────────────────────────────────
\subsection{Structure}

- **Data**: 1,000 hourly response time observations (synthetic)
- **Regime 1** (obs 0--599): baseline mean 200ms, std 30ms, slight daily seasonality
- **Regime 2** (obs 600--999): elevated mean 350ms, std 30ms (infrastructure degradation)
- **Anomalies**: 15 injected spike anomalies at random positions (5x normal magnitude)
- **Questions**: (1) Which observations are anomalous spikes? (2) Where does the regime change occur?

% ── Mapping ───────────────────────────────────────────────────────────
\subsection{Real-World Mapping}

\begin{longtable}{ll}
\toprule
\textbf{Real-World Concept} & \textbf{Mathematical Object} \\
\midrule
\endhead
problem input & $instance parameters$ \\
\bottomrule
\end{longtable}

% ── Constraints ───────────────────────────────────────────────────────
\subsection{Constraints}

\begin{enumerate}
  \item $Anomaly detection**: identifying observations that deviate significantly from expected behavior$
  \hfill \textit{()}
  \item $Change point detection**: finding the time index where the statistical properties of a series shift$
  \hfill \textit{()}
  \item $Z-score**: number of standard deviations from the (rolling) mean; threshold-based flagging$
  \hfill \textit{()}
  \item $IQR (interquartile range)**: nonparametric spread measure; outliers fall beyond 1.5*IQR fences$
  \hfill \textit{()}
  \item $PELT**: Pruned Exact Linear Time algorithm for optimal segmentation (Killick et al. 2012)$
  \hfill \textit{()}
  \item $Rolling statistics**: windowed mean/std that adapt to local behavior, preventing regime shifts from inflating global statistics$
  \hfill \textit{()}
  \item $Precision / recall / F1**: standard classification metrics applied to anomaly detection performance$
  \hfill \textit{()}
\end{enumerate}

% ── Objective / Claim ─────────────────────────────────────────────────
\subsection{Claim}


% ── Approach ──────────────────────────────────────────────────────────
\subsection{Approach}

\begin{description}
  \item[Suggested approach:] Z-score (rolling) + IQR + PELT (L2, BIC penalty)
  \item[Complexity class:] 
  \item[Available tools:] Z-score
\end{description}
% ── Phase B: Solution ────────────────────────────────────────────────
\section{Solution}

\begin{polyaresult}

\textbf{Answer}: Solution computed


\textbf{Optimal}: No / Unknown \hfill
\textbf{Feasible}: No
\textbf{Algorithm}: Z-score (rolling) + IQR + PELT (L2, BIC penalty) \quad $$

\textbf{Time}: 3.1975480560213327s


\end{polyaresult}

% ── Solution Details ──────────────────────────────────────────────────
\subsection{Solution Details}


Method: 1. Z-score (rolling window): Compute rolling mean and std (window=50). Flag observations where |z-score| > 3.0. Adapts to local trends and is robust to gradual shifts.

2. IQR method (rolling window): For each observation, compute Q1, Q3, IQR over a local window (100 points). Flag observations beyond Q1 - 1.5*IQR or Q3 + 1.5*IQR fences. The rolling approach adapts to regime changes, unlike a global IQR. Simple, nonparametric, resistant to distributional assumptions.

3. PELT change point detection (ruptures library): Pruned Exact Linear Time algorithm with L2 (least squares) cost model and BIC penalty. Finds the optimal segmentation of the time series into regimes with different means.

% ── Verification ──────────────────────────────────────────────────────
\subsection{Verification}

\begin{longtable}{lcl}
\toprule
\textbf{Check} & \textbf{Status} & \textbf{Detail} \\
\midrule
\endhead
Recall Gte 0.7 & \passmark & Passed \\
Precision Gte 0.5 & \passmark & Passed \\
Change Point Within 20 & \passmark & Passed \\
Few Change Points & \passmark & Passed \\
Regime Diff Significant & \passmark & Passed \\
Zscore Anomalies Extreme & \passmark & Passed \\
All Passed & \passmark & Passed \\
\bottomrule
\end{longtable}

% ── Phase C: Interpretation ──────────────────────────────────────────
\section{Interpretation}

% ── The Question & The Answer ─────────────────────────────────────────
\subsection{The Question}
- **Data**: 1,000 hourly response time observations (synthetic)
- **Regime 1** (obs 0--599): baseline mean 200ms, std 30ms, slight daily seasonality
- **Regime 2** (obs 600--999): elevated mean 350ms, std 30ms (infrastructure degradation)
- **Anomalies**: 15 injected spike anomalies at random positions (5x normal magnitude)
- **Questions**: (1) Which observations are anomalous spikes? (2) Where does the regime change occur?.

\subsection{The Answer}

\begin{polyainsight}[title={Bottom Line}]
A feasible solution was found.
\end{polyainsight}

% ── What This Means ───────────────────────────────────────────────────
\subsection{What This Means}

- Z-score anomaly detection: detects most injected spikes (recall >= 70\%, precision >= 50\%)
- IQR anomaly detection: complementary method with similar recall
- PELT change point: detected within 20 observations of the true change at obs 600
- Pre/post mean difference: statistically significant (t-test p < 0.001)
- All 6 verification checks pass

1. Z-score (rolling window): Compute rolling mean and std (window=50). Flag observations where |z-score| > 3.0. Adapts to local trends and is robust to gradual shifts.

2. IQR method (rolling window): For each observation, compute Q1, Q3, IQR over a local window (100 points). Flag observations beyond Q1 - 1.5*IQR or Q3 + 1.5*IQR fences. The rolling approach adapts to regime changes, unlike a global IQR. Simple, nonparametric, resistant to distributional assumptions.

3. PELT change point detection (ruptures library): Pruned Exact Linear Time algorithm with L2 (least squares) cost model and BIC penalty. Finds the optimal segmentation of the time series into regimes with different means.

% ── Sensitivity Analysis ──────────────────────────────────────────────

% ── Figures ───────────────────────────────────────────────────────────

% ── Recommendations ───────────────────────────────────────────────────
\subsection{Recommendations}

\begin{enumerate}
  \item Review the model assumptions to ensure real-world fidelity.
\end{enumerate}

% ── Limitations ───────────────────────────────────────────────────────
\subsection{Limitations}

\begin{polyawarnbox}
\begin{itemize}
  \item Model assumes deterministic parameters; real-world variability not captured.
  \item Solution quality depends on the accuracy of input data.
\end{itemize}
\end{polyawarnbox}

% ── Appendix ─────────────────────────────────────────────────────────

\end{document}